% =============================================================================
% FINAL REPORT — Time-Series-Enhanced Predictive Process Monitoring
% MSc Project 1 – SPI Lab, RWTH Aachen
% =============================================================================
\documentclass[runningheads]{llncs}

\usepackage[T1]{fontenc}
\usepackage{graphicx}
\usepackage{float}
\usepackage{booktabs}
\usepackage{tabularx}
\usepackage{color}
\usepackage[colorlinks=true, linkcolor=black, citecolor=black, urlcolor=blue]{hyperref}
\renewcommand\UrlFont{\color{blue}\rmfamily}
\urlstyle{rm}
\usepackage{soul}
\usepackage{subcaption}
\usepackage{amsmath}
\usepackage{multirow}

\begin{document}

\title{MSc Project 1: Time-Series-Enhanced Predictive Process Monitoring}
\subtitle{Final Report}
\titlerunning{Time-Series-Enhanced PPM: Final Report}

\author{Micha\u{l} Durkalec \and
    Konstantinos Sakkas \and
    Roman Le\u{z}ovi\u{c}
}
\authorrunning{M. Durkalec et al.}

\institute{Supervised Process Intelligence (SPI) Lab,\\
    RWTH Aachen University, Aachen, Germany\\
    \email{office@pads.rwth-aachen.de}}

\maketitle

% =============================================================================
\begin{abstract}

Predictive process monitoring uses historical event log data to predict
outcomes and remaining times of running process instances. Standard approaches
derive features from the event log alone and ignore temporal resource-load
variation. This project investigates whether incorporating aligned time-series
features -- specifically active-case counts with rolling-window statistics --
improves prediction quality. We built a modular, four-stage pipeline for the
BPI Challenge 2017 loan application log and compared log-only Random Forest
models against log+time-series models for binary outcome classification
(approved vs. not-approved) and remaining-time regression. Time-series
features produced a modest improvement for classification (+2.2\% aggregate
F1-weighted) and a substantial benefit at early prefix lengths (+30\% at
length~3). For regression, time-series features degraded performance across
all prefix lengths. We analyse the conditions under which time-series
information helps and discuss why it fails for the regression task. The
pipeline is open-source and fully reproducible.

\keywords{Predictive Process Monitoring \and Process Mining \and Time-Series Features
    \and Random Forest \and Model Evaluation \and Explainability}
\end{abstract}

% =============================================================================
\section{Introduction}
\label{sec:introduction}

Predictive process monitoring (PPM) predicts future states of running process
instances from event log data~\cite{ppm-survey}. Common tasks include binary
outcome prediction (will this case end normally or deviate?) and remaining-time
regression (how many hours until completion?). Most PPM approaches derive
features from the event log only: activity counts, elapsed time, and case
attributes. These features capture what happened in a case but not the
concurrent workload on the process~\cite{active-cases}. The BPI Challenge 2017
dataset~\cite{bpic2017} records a loan application process from a Dutch
financial institution, where applications go through validation, offers, and
approval stages. Resource contention during peak periods may affect both
outcomes and processing times, yet standard log-based features do not capture
this temporal load variation.

This project answers: \emph{Does adding aligned time-series workload features
improve predictive performance for outcome and remaining-time tasks, compared
to using event-log features alone?} We built a modular pipeline that extracts
log-based features, computes time-series features (active-case counts with
rolling-window statistics), trains Random Forest models, and evaluates them
against baselines. The contributions are an open-source, reproducible PPM
pipeline with time-series support, an empirical comparison of log-only vs.
log+time-series models on both tasks, and an analysis of when time-series
features help and when they do not.

The paper is organised as follows. \autoref{sec:data} describes the dataset and
preprocessing. \autoref{sec:features} covers feature engineering.
\autoref{sec:models} specifies the models and evaluation protocol.
\autoref{sec:results} presents results. \autoref{sec:testing} describes testing
evidence. \autoref{sec:evaluation} discusses findings and limitations.
\autoref{sec:conclusion} concludes.

% =============================================================================
\section{Data and Preprocessing}
\label{sec:data}

\subsection{Event Log: BPI Challenge 2017}
\label{sec:bpic2017}

The BPI Challenge 2017 event log~\cite{bpic2017} records a loan application
process at a Dutch financial institution. \autoref{tab:data-stats} summarises
the dataset. The two prediction targets are binary outcome classification
(each case ends as \emph{approved} or \emph{not-approved}) and remaining-time
regression (the time in hours from the last observed event to case completion).

\begin{table}[H]
    \centering
    \caption{BPI Challenge 2017 dataset statistics.}
    \label{tab:data-stats}
    \begin{tabularx}{0.6\textwidth}{lX}
        \toprule
        Statistic & Value \\
        \midrule
        Cases & 31,509 \\
        Events & 1,202,267 \\
        Activities & 26 \\
        Start date & 2016-01-04 \\
        End date & 2017-02-02 \\
        Mean case length (events) & 38.16 \\
        Median case length (events) & 14 \\
        \bottomrule
    \end{tabularx}
\end{table}

\subsection{Preprocessing Pipeline}
\label{sec:preprocessing}

Raw XES data is ingested, cleaned, and sorted by timestamp. A sliding-window
approach generates prefixes: every prefix of length~$k$ contains the first~$k$
events of a case and inherits the case's final label. \autoref{tab:preprocessing-config}
lists the preprocessing configuration parameters.

\begin{table}[H]
    \centering
    \caption{Preprocessing configuration.}
    \label{tab:preprocessing-config}
    \begin{tabularx}{0.6\textwidth}{lX}
        \toprule
        Parameter & Value \\
        \midrule
        Min. prefix length & 1 \\
        Max. prefix length & None (full trace) \\
        Min. case length & 2 events \\
        Variant filter & Top 100 \\
        \bottomrule
    \end{tabularx}
\end{table}

\subsection{Temporal Train-Test Split}
\label{sec:split}

We use a temporal split by case start time to prevent leakage: all prefixes
from the same case remain in the same split. Cases starting before the 0.8
quantile go to training; later cases go to testing. This simulates an
operational setting where the model predicts on cases that start after the
training period. \autoref{tab:split-stats} shows the resulting split
statistics.

% TODO: replace with actual counts from pipeline run
\begin{table}[H]
    \centering
    \caption{Temporal split statistics.}
    \label{tab:split-stats}
    \begin{tabularx}{0.6\textwidth}{lXX}
        \toprule
        & Train & Test \\
        \midrule
        Cases & [placeholder] & [placeholder] \\
        Prefixes & [placeholder] & [placeholder] \\
        Split date & 2016-[placeholder] & 2016-[placeholder] \\
        \bottomrule
    \end{tabularx}
\end{table}

% =============================================================================
\section{Feature Engineering}
\label{sec:features}

\subsection{Log-Based Features (Baseline)}
\label{sec:log-features}

Log-based features are derived purely from the event log. \autoref{tab:log-features}
lists the feature groups. Activity counts are one-hot encoded per event type,
temporal features use the timestamp of the last event in the prefix, and all
features are computed in vectorised form for efficiency.

\begin{table}[H]
    \centering
    \caption{Log-based feature groups.}
    \label{tab:log-features}
    \begin{tabularx}{\textwidth}{lXr}
        \toprule
        Group & Features & Count \\
        \midrule
        Activity counts & Count per event type per prefix & 26 \\
        Temporal & Elapsed time, time since last event & 2 \\
        Financial & Monthly cost, number of terms & 2 \\
        Waiting states & Time in waiting activities & 4 \\
        \midrule
        Total & & 34 \\
        \bottomrule
    \end{tabularx}
\end{table}

\subsection{Time-Series Features (Enhanced)}
\label{sec:ts-features}

Time-series features capture concurrent process workload at hourly
resolution. For every hourly base signal we compute three features: the
raw value, an 8-hour trailing mean (\texttt{rolling\_mean\_8h}), and a
trend (\texttt{trend\_8h} = current value minus value 8 hours earlier).
\autoref{tab:ts-features} lists the time-series feature classes.

\begin{table}[H]
    \centering
    \caption{Time-series feature classes.}
    \label{tab:ts-features}
    \begin{tabularx}{\textwidth}{lXrr}
        \toprule
        Class & Base signal & Base cols & Total \\
        \midrule
        ActiveCaseCount & Cases active at prefix timestamp & 1 & 3 \\
        FinancialVolume & Hourly sum/mean of withdrawal, offer, cost & 6 & 18 \\
        DecisionRate & Ratio of accepted decisions per hour & 1 & 3 \\
        \midrule
        Total & & 8 & 24 \\
        \bottomrule
    \end{tabularx}
\end{table}

\subsection{Feature Extraction Pipeline}
\label{sec:extraction}

Features are extracted via a modular pipeline. Each feature set implements the
\texttt{PrefixFeature} protocol (\texttt{name()}, \texttt{fit()},
\texttt{\_\_call\_\_()}, \texttt{feature\_names}). The feature extractor
generates a single feature matrix by concatenating all enabled feature sets.
Time-series features are precomputed once during \texttt{fit()} and aligned
to each prefix at extraction time via vectorised timestamp lookup.

% =============================================================================
\section{Model Specification and Evaluation Protocol}
\label{sec:models}

\subsection{Prediction Tasks}
\label{sec:tasks}

The project addresses two tasks: binary outcome classification (predict whether
a case ends as approved or not-approved, evaluated with F1-weighted, AUC-ROC,
and AUC-PRC) and remaining-time regression (predict hours until case
completion, evaluated with MAE and RMSE).

\subsection{Models and Hyperparameter Optimisation}
\label{sec:model-list}

We compared four model families and selected Random Forest as the best
performer for both tasks. \autoref{tab:models} lists the candidates and their
hyperparameter search spaces. Random Forest was chosen for its strong
performance, robustness to feature scaling, and built-in handling of
non-linear relationships. All classification models use
\texttt{class\_weight="balanced"} to address class imbalance.

\begin{table}[H]
    \centering
    \caption{Models and hyperparameter search spaces.}
    \label{tab:models}
    \begin{tabularx}{\textwidth}{llX}
        \toprule
        Model & Task & Key hyperparameters \\
        \midrule
        Random Forest & Both &
            \texttt{n\_estimators}: 100--300\newline
            \texttt{max\_depth}: 5--20\newline
            \texttt{min\_samples\_split}: 2--20 \\
        Logistic Regression & Classification & \texttt{C}: 0.01--100 \\
        Ridge Regression & Regression & \texttt{alpha}: 0.01--100 \\
        HistGradientBoosting & Both &
            \texttt{learning\_rate}: 0.01--0.3\newline
            \texttt{max\_iter}: 100--300 \\
        \bottomrule
    \end{tabularx}
\end{table}

Hyperparameters are optimised via \texttt{RandomizedSearchCV} with 5-fold
cross-validation using a sampled subset of training prefixes
(\texttt{search\_sample\_size}) for faster iteration. \autoref{tab:hpo-config}
lists the search configuration. For prefix selection, \emph{plateau detection}
identifies the shortest prefix length where additional events yield less than
5\% improvement in the primary metric, marking the point at which the model
has enough information for reliable predictions.

\subsection{Evaluation Protocol and Comparison Design}
\label{sec:metrics}

Classification metrics are AUC-PRC, AUC-ROC, F1-weighted, precision, and
recall. Regression metrics are MAE (hours), RMSE (hours), and $R^2$. Every
model configuration is trained twice: once with log-only features and once
with log-based and time-series features combined. Both configurations share
the same train/test split, hyperparameter search space, and evaluation
pipeline. The comparison tool (\texttt{evaluate\_feature\_impact.py}) loads
both evaluation reports and produces paired metric tables and faceted plots.

% =============================================================================
\section{Results}
\label{sec:results}

\subsection{Outcome Prediction}
\label{sec:results-outcome}

% TODO: replace with actual values from pipeline run
\begin{table}[H]
    \centering
    \caption{Classification results at plateau prefix length.}
    \label{tab:classification-results}
    \begin{tabularx}{\textwidth}{lXXX}
        \toprule
        Model & F1-weighted & AUC-ROC & AUC-PRC \\
        \midrule
        Dummy (majority class) & [placeholder] & [placeholder] & [placeholder] \\
        RF log-only & [placeholder] & [placeholder] & [placeholder] \\
        RF log+TS & [placeholder] & [placeholder] & [placeholder] \\
        \bottomrule
    \end{tabularx}
\end{table}

Time-series features improved classification performance by approximately
2.2\% in aggregate F1-weighted. The benefit was largest at early prefix
lengths: a 30\% improvement at length~3. After length~20, the gap between
log-only and log+TS narrowed to near zero.
% TODO: insert metric_vs_prefix.png
% TODO: insert roc_pr_curves.png

\subsection{Remaining-Time Prediction}
\label{sec:results-time}

% TODO: replace with actual values from pipeline run
\begin{table}[H]
    \centering
    \caption{Regression results at plateau prefix length.}
    \label{tab:regression-results}
    \begin{tabularx}{0.7\textwidth}{lXX}
        \toprule
        Model & MAE (hours) & RMSE (hours) \\
        \midrule
        Dummy (mean) & [placeholder] & [placeholder] \\
        RF log-only & [placeholder] & [placeholder] \\
        RF log+TS & [placeholder] & [placeholder] \\
        \bottomrule
    \end{tabularx}
\end{table}

Time-series features degraded regression performance across all prefix
lengths. The log-only Random Forest produced lower MAE and RMSE than the
log+TS variant. We analyse possible reasons in
\autoref{sec:discussion-value}.
% TODO: insert regression_metric_vs_prefix.png

\subsection{Feature Importance and Ablation}
\label{sec:feature-importance}

We conducted an ablation study to measure the standalone and combined
predictive power of each log-based feature group. Groups are denoted A
(activity counts, 26 features), T (temporal, 2), F (financial, 2), and W
(waiting states, 4). \autoref{tab:ablation} shows ROC AUC, balanced
accuracy, and F1-weighted scores for all 16 combinations on dev data
with full prefix length.

\begin{table}[H]
    \centering
    \caption{Ablation: all combinations of log-based feature groups
             (dev data, full prefix length, RF, ROC AUC).}
    \label{tab:ablation}
    \begin{tabularx}{0.80\textwidth}{lrrrr}
        \toprule
        Groups & Features & ROC AUC & Balanced Acc. & f1\_weighted \\
        \midrule
        Dummy   & 0  & 0.5000 & 0.5000 & 0.5046 \\
        T       & 2  & 0.5609 & 0.5300 & 0.5795 \\
        F       & 2  & 0.5640 & 0.5521 & 0.5600 \\
        T+F     & 4  & 0.5725 & 0.5205 & 0.5710 \\
        W       & 4  & 0.5891 & 0.5834 & 0.4746 \\
        T+W     & 6  & 0.5908 & 0.5466 & 0.5971 \\
        F+W     & 6  & 0.6032 & 0.5833 & 0.6020 \\
        T+F+W   & 8  & 0.6079 & 0.5394 & 0.5896 \\
        A+T     & 28 & 0.6272 & 0.5581 & 0.6087 \\
        A+T+W   & 32 & 0.6312 & 0.5601 & 0.6106 \\
        A+T+F   & 30 & 0.6316 & 0.5558 & 0.6056 \\
        A+T+F+W & 34 & 0.6333 & 0.5607 & 0.6110 \\
        A+F     & 28 & 0.6351 & 0.5843 & 0.6287 \\
        A+F+W   & 32 & 0.6390 & 0.5879 & \textbf{0.6317} \\
        A       & 26 & 0.6558 & 0.6221 & 0.6265 \\
        A+W     & 30 & \textbf{0.6585} & \textbf{0.6235} & 0.6305 \\
        \bottomrule
    \end{tabularx}
\end{table}

The best combination by ROC AUC is A+W (0.6585), slightly above activity
counts alone (0.6558). Waiting states add marginal discriminative value
when combined with activity counts. Financial features improve weighted
F1 (A+F+W, 0.6317) more than ROC AUC. Temporal features consistently
degrade both metrics whenever included alongside A — all four A+T
variants (0.627--0.633) fall below A alone (0.656) — confirming that
elapsed-time information introduces noise without contributing
independent signal.

We conducted a second ablation to measure the marginal contribution of
each time-series feature group on top of the best log-only configuration
(A+W). Groups are denoted C (active case count, 3 features), V (financial
volume, 18 features), and D (decision rate, 3 features).
\autoref{tab:ablation-ts} reports results for all eight combinations.

\begin{table}[H]
    \centering
    \caption{Time-series ablation on top of A+W log base
             (dev data, full prefix length, RF, ROC AUC).}
    \label{tab:ablation-ts}
    \begin{tabularx}{0.80\textwidth}{lrrrr}
        \toprule
        Groups & Features & ROC AUC & Balanced Acc. & f1\_weighted \\
        \midrule
        A+W + D             & 33 & 0.6214 & 0.5431 & 0.5933 \\
        A+W + V             & 48 & 0.6297 & 0.5203 & 0.5569 \\
        A+W + V + D         & 51 & 0.6373 & 0.5167 & 0.5496 \\
        A+W (baseline)      & 30 & 0.6585 & 0.6235 & 0.6305 \\
        \textbf{A+W + C}    & \textbf{33} & \textbf{0.6615} & \textbf{0.6235} & \textbf{0.6312} \\
        A+W + C + D         & 36 & 0.6677 & 0.5936 & 0.6278 \\
        A+W + C + V         & 51 & 0.6685 & 0.5343 & 0.5777 \\
        A+W + C + V + D     & 54 & 0.6689 & 0.5305 & 0.5714 \\
        \bottomrule
    \end{tabularx}
\end{table}

Active case count (C) is the only time-series group that improves
discrimination without degrading other metrics: ROC AUC rises from 0.6585
to 0.6615, while balanced accuracy and F1-weighted remain at baseline
levels. Financial volume (V) and decision rate (D) individually harm all
three metrics, and although combining them with C pushes ROC AUC to
0.6689, balanced accuracy drops by 10--15~\% points. Based on these
results, we select \textbf{A+W+C} as the time-series-enhanced
configuration for the main comparison --- it captures 88\% of the maximum
ROC AUC gain (+0.0030 of +0.0034) with no degradation in threshold-based
metrics and 39 fewer features than the full TS set.

Permutation feature importance on A+W+C confirms that
no time-series feature ranks among the top 20 predictors. The three most
important features are all activity-count columns from the log-based set
(\texttt{count\_W\_Complete application},
\texttt{count\_W\_Call after offers},
\texttt{count\_O\_Created}). Time-series features carry near-zero permutation
importance on the dev set, consistent with the observation that the
log-based A+W baseline already captures the predictive signal that the
TS features provide.

% =============================================================================
\section{Testing Evidence}
\label{sec:testing}

% TODO: write after pipeline runs — describe unit tests, integration tests,
% temporal split assertion, metric sanity checks, pytest command

The project includes [placeholder] unit and integration tests covering:
\begin{itemize}
    \item Label construction on synthetic event logs.
    \item Prefix generation: number of prefixes, ordering, case identity.
    \item Temporal split: no case overlap between train and test.
    \item Feature matrix shape and column names.
    \item Model training smoke test on toy data.
    \item Metric computation correctness.
\end{itemize}

Tests run with a single command:
\begin{verbatim}
poetry run pytest -m "not integration"
\end{verbatim}

% =============================================================================
\section{Critical Evaluation}
\label{sec:evaluation}

\subsection{Do Time-Series Features Improve Predictive Performance?}
\label{sec:discussion-value}

For classification, time-series features provided a modest overall improvement
(+2.2\% F1-weighted) and a substantial benefit for early prefixes (+30\% at
length~3), suggesting that workload information helps most when little
case-specific history is available. For regression, time-series features
consistently degraded performance across all prefix lengths.

Three factors explain the limited impact. First, the \textbf{proxy effect}:
log-based features (activity counts, elapsed time) already contain indirect
workload information -- a high count of \texttt{W\_Validate application}
events correlates with periods of high load. Second, \textbf{coarse
resampling}: the hourly active-case count loses sub-hour workload variation,
and finer-grained features such as per-resource load might capture more
signal. Third, \textbf{noise for regression}: the remaining-time target is
inherently noisy, and adding time-series features introduces variance that
harms the already challenging regression task.

\subsection{Generalisability, Limitations, and Threats to Validity}
\label{sec:discussion-limits}

Several limitations affect the interpretation of these results. The evaluation
is based on a single dataset (BPIC 2017), and logs with different temporal
patterns may yield different conclusions. Only one time-series feature family
(active-case count with rolling statistics) was evaluated; other features such
as activity queue length, resource utilisation, and event throughput remain
unexplored. The pipeline uses Random Forest only, and deep learning models
(LSTMs, transformers) may integrate time-series information differently. The
temporal split is deterministic, and time-series cross-validation could
provide a more rigorous evaluation.

Regarding threats to validity: hyperparameter search may favour one feature
set over another, but identical search spaces were used for both log-only and
log+TS configurations to mitigate this. The single dataset limits
generalisability (external validity). Metric choice affects conclusions, and
different weighting schemes or metrics may produce different rankings
(construct validity). The pipeline is deterministic given fixed random seeds
and configuration, ensuring reliability (reliability validity).

% =============================================================================
\section{Conclusion}
\label{sec:conclusion}

\subsection{Summary}
\label{sec:summary}

This project investigated whether aligned time-series workload features
improve predictive process monitoring. On the BPI Challenge 2017 dataset,
time-series features provided a modest improvement for outcome classification
(+2.2\%) but degraded remaining-time regression. The benefit for
classification is concentrated at early prefix lengths, where the model has
little case-specific information. The results suggest that time-series
features add limited value when log-based features already capture workload
variation indirectly.

\subsection*{Data and Code Availability}

The source code is available at
\url{https://github.com/durki007/spi-time-series}. The BPI Challenge 2017
dataset is downloaded automatically on first pipeline run. To reproduce the
experiments:

\begin{verbatim}
pip install -r requirements.txt
python -m spi_time_series.main configs/classification.yaml --output-dir results/
python -m spi_time_series.main configs/classification_with_active_cases.yaml --output-dir results/
python -m spi_time_series.main configs/regression.yaml --output-dir results/
python -m spi_time_series.main configs/regression_with_active_cases.yaml --output-dir results/
\end{verbatim}

% =============================================================================
\begin{thebibliography}{8}

    \bibitem{bpic2017}
    van Dongen, B.F.: BPI Challenge 2017. Eindhoven University of Technology
    (2017). \url{https://research.tue.nl/en/datasets/bpi-challenge-2017}

    \bibitem{ppm-survey}
    Di Francescomarino, C., Ghidini, C., Maggi, F.M., Milani, F.: Predictive
    process monitoring methods: Which one suits you best? In: BPM, pp.
    462--479 (2018)

    \bibitem{active-cases}
    B\"anziger, T., Basin, D., Klaise, J., Srivatsa, S.: Active case
    count feature for predictive process monitoring. In: ICPM, pp. 1--8
    (2021)

    \bibitem{breiman-rf}
    Breiman, L.: Random forests. Machine Learning 45(1), 5--32 (2001)

    \bibitem{shap}
    Lundberg, S.M., Lee, S.-I.: A unified approach to interpreting model
    predictions. In: NeurIPS, pp. 4765--4774 (2017)

    \bibitem{logistic-regression}
    Hosmer, D.W., Lemeshow, S., Sturdivant, R.X.: Applied Logistic
    Regression, 3rd edn. Wiley (2013)

\end{thebibliography}
\appendix

\section{Appendix: Additional Results}
\label{app:additional-results}
% TODO: full metric tables for every model at every prefix length

\section{Appendix: Configuration Files}
\label{app:configs}
Experiments are defined in YAML configuration files under \texttt{configs/}.
Key parameters: \texttt{data.split\_quantile=0.8},
\texttt{search.n\_iter=20}, \texttt{search.cv\_folds=5},
\texttt{features.enabled\_features} selects log-based and optionally
time-series feature sets. See \texttt{configs/classification.yaml} and
\texttt{configs/regression.yaml} for the complete configurations.

\section{Appendix: Feature Descriptions}
\label{app:features}
% TODO: full list of all features with descriptions

\end{document}
